\chapter{Correspondência de notação}
\label{cap:notacao}

A notação do livro foi harmonizada. Este apêndice registra a correspondência com os
símbolos dos trabalhos originais, para que o leitor possa transitar entre as fontes sem
ambiguidade. A tabela de símbolos propriamente dita encontra-se nas páginas pré-textuais.

\section{Séries temporais e janelas}

\begin{table}[htbp]
\centering
\caption{Correspondência de notação para dados e janelas.}
\label{tab:notb-dados}
\small
\begin{tabular}{p{2.3cm}p{4.6cm}p{5.4cm}}
\toprule
\textbf{Neste livro} & \textbf{Significado} & \textbf{Nos trabalhos originais} \\
\midrule
$T$ &
Passos de tempo da janela &
\emph{window size}, \emph{sequence length}, $n_{\text{steps}}$; valores 200, 500 e a
varredura de 150 a 7500 do \cref{cap:closedworld} \\
\addlinespace
$S$ &
Número de sensores &
\emph{n\_features}, \emph{channels}; valores 4 e 5 conforme o capítulo \\
\addlinespace
$\mat{X} \in \R^{T \times S}$ &
Janela multivariada &
$X$, \emph{input tensor}, com forma $(T, S)$ na notação de \ing{frameworks} \\
\addlinespace
$\vect{b}$, $\vect{e}$ &
Janelas de linha de base e de evento &
\emph{baseline window} e \emph{event window} no \cref{cap:agente} \\
\addlinespace
$y$ &
Rótulo de classe &
\emph{class}, \emph{label}; no 3W, o código de 0 a 9 \\
\bottomrule
\end{tabular}
\end{table}

\begin{destaque}[Uma diferença de convenção que confunde]
Bibliotecas de aprendizado profundo descrevem a entrada como \emph{(batch, timesteps,
features)}. Este livro omite a dimensão de lote e escreve $\mat{X} \in \R^{T \times S}$.

Assim, a camada descrita nos trabalhos originais como \emph{Input(500, 5)} corresponde
aqui a $T = 500$ e $S = 5$.
\end{destaque}

\section{Modelos e representações}

\begin{table}[htbp]
\centering
\caption{Correspondência de notação para modelos.}
\label{tab:notb-modelos}
\small
\begin{tabular}{p{2.3cm}p{4.6cm}p{5.4cm}}
\toprule
\textbf{Neste livro} & \textbf{Significado} & \textbf{Nos trabalhos originais} \\
\midrule
$g_{\phi}$, $h_{\psi}$ &
Codificador e decodificador &
\emph{encoder} e \emph{decoder}; frequentemente sem símbolo, apenas nomeados \\
\addlinespace
$\vect{z}$ &
Representação latente &
\emph{latent vector}, \emph{embedding}, \emph{bottleneck}, \emph{latent
representation} \\
\addlinespace
$d$ &
Dimensão do latente &
\emph{latent dim}; valores 4, 16, 64 e 128 conforme o capítulo \\
\addlinespace
$c_k$, $s_k$ &
Classificador binário da classe $k$ e seu escore &
\emph{binary classifier k}, \emph{sigmoid output}, \emph{score} \\
\addlinespace
$\tau$ &
Limiar de decisão &
\emph{threshold}; valores 0,3, 0,5 e 0,65 conforme o contexto \\
\bottomrule
\end{tabular}
\end{table}

\begin{destaque}[Três latentes distintos, um mesmo símbolo]
O símbolo $\vect{z}$ designa, em capítulos diferentes, objetos com propriedades muito
distintas:

no \cref{cap:mundoaberto}, o latente de dimensão 16 de um autoencoder;

no \cref{cap:mahalanobis}, o latente de dimensão 128 de um autoencoder maior \emph{e} a
penúltima camada de dimensão 64 de um classificador multiclasse.

A tese central do \cref{cap:mahalanobis} é justamente que esses espaços não são
intercambiáveis, apesar de receberem o mesmo símbolo. Quando a distinção importa, o texto
a explicita.
\end{destaque}

\section{Estatística e limiares}

\begin{table}[htbp]
\centering
\caption{Correspondência de notação para estatística e limiares.}
\label{tab:notb-estat}
\small
\begin{tabular}{p{2.3cm}p{4.6cm}p{5.4cm}}
\toprule
\textbf{Neste livro} & \textbf{Significado} & \textbf{Nos trabalhos originais} \\
\midrule
$\mu + k\sigma$ &
Limiar de reconstrução &
\emph{mean + 3 std}; no \cref{cap:dual}, $k = 3$ \\
\addlinespace
$d_M$ &
Distância de Mahalanobis &
\emph{Mahalanobis distance}, $D^2$ quando elevada ao quadrado \\
\addlinespace
$\lambda$ &
Regularização da covariância &
\emph{shrinkage}, \emph{regularization term} somado à diagonal \\
\addlinespace
$K$ &
Número de agrupamentos &
$k$ minúsculo em \emph{k-means}; aqui maiúsculo, para evitar colisão \\
\addlinespace
$R = \mu/\sigma$ &
Razão do limiar adaptativo &
\emph{ratio} no \cref{cap:mundoaberto} \\
\addlinespace
$\delta = \bar{e} - \bar{b}$ &
Variação absoluta &
\emph{delta}, \emph{absolute change} no \cref{cap:agente} \\
\bottomrule
\end{tabular}
\end{table}

\section{Métricas}

\begin{table}[htbp]
\centering
\caption{Correspondência de notação para métricas de avaliação.}
\label{tab:notb-metricas}
\small
\begin{tabular}{p{2.3cm}p{4.6cm}p{5.4cm}}
\toprule
\textbf{Neste livro} & \textbf{Significado} & \textbf{Nos trabalhos originais} \\
\midrule
ACC & Acurácia & \emph{accuracy}, \emph{acc} \\
$\text{ACC}_b$ & Acurácia balanceada & \emph{balanced accuracy} \\
REC & Revocação & \emph{recall}, \emph{sensitivity}, \emph{true positive rate} \\
PR & Precisão & \emph{precision}, \emph{positive predictive value} \\
SP & Especificidade & \emph{specificity}, \emph{true negative rate} \\
$F_1$ & Medida F1 & \emph{F1-score}, \emph{F-measure} \\
IoU & Interseção sobre união & \emph{IoU}, \emph{Jaccard index} \\
\bottomrule
\end{tabular}
\end{table}

\begin{destaque}[Duas armadilhas de tradução]
\textbf{``Precisão''} traduz \emph{precision} neste livro, e nunca \emph{accuracy}. Em
português técnico corrente, ``precisão'' é usada às vezes para acurácia, o que gera
confusão. Aqui, acurácia é ACC e precisão é PR.

\textbf{``Sensibilidade''} e \textbf{``revocação''} são o mesmo objeto. O livro usa
revocação uniformemente, mas a literatura médica e de sinais prefere sensibilidade.
\end{destaque}

\section{Nomenclatura de classes}

\begin{table}[htbp]
\centering
\caption{Correspondência entre a numeração de classes do 3W e a designação usada no
livro.}
\label{tab:notb-classes}
\small
\begin{tabular}{clp{6.8cm}}
\toprule
\textbf{Código} & \textbf{Designação no livro} & \textbf{Denominação no 3W} \\
\midrule
0 & Normal & \emph{Normal} \\
1 & Aumento de BSW & \emph{Abrupt Increase of BSW} \\
2 & Fechamento de DHSV & \emph{Spurious Closure of DHSV} \\
3 & Golfadas severas & \emph{Severe Slugging} \\
4 & Instabilidade de fluxo & \emph{Flow Instability} \\
5 & Perda de produtividade & \emph{Rapid Productivity Loss} \\
6 & Restrição rápida no PCK & \emph{Quick Restriction in PCK} \\
7 & Incrustação no PCK & \emph{Scaling in PCK} \\
8 & Hidrato na produção & \emph{Hydrate in Production Line} \\
9 & Hidrato no serviço & \emph{Hydrate in Service Line} \\
\bottomrule
\end{tabular}
\end{table}

\begin{destaque}[Deslocamento de índice]
Alguns dos trabalhos originais indexam as nove classes anômalas de 1 a 9; outros indexam
de 0 a 8, tratando a classe normal separadamente. Este livro segue estritamente a
numeração do 3W, de 0 a 9, com a classe 0 correspondendo ao regime normal.

Ao comparar tabelas com as fontes originais, verifique a convenção antes de alinhar
linhas.
\end{destaque}
